\subsection{A norm bound for hidden alternatives}

Exact fixing quantifies over every admissible difference, however small.
To ask about its size, a restriction on alternatives must be supplied. Use
the conformal Cauchy-data norm
\[
 N(f,v)^2=\|\Lambda f\|_2^2+\|v\|_2^2
\]
at the auxiliary slice $\eta=0$, with unit-$S^3$ volume $2\pi^2$.
This is not an observer's measured physical-energy budget.

\begin{proposition}\label{prop:norm}
Let $\psi_\lambda$ and $\psi_0$ have the same smooth data $(f,v)$ and solve
the cylinder equation with spatially homogeneous smooth real potential
$\lambda q(\eta)$ and with zero potential, respectively. For $\eta\ge0$, put
$Q(\eta)=\int_0^\eta |q(u)|\,du$. Then
\begin{equation}\label{eq:norm}
 \sup_X|\psi_\lambda(\eta,X)-\psi_0(\eta,X)|
 \le \frac{|\lambda|Q(\eta)e^{|\lambda|Q(\eta)}}{\sqrt6}\,N(f,v).
\end{equation}
If $f=0$, the denominator improves to $\sqrt{12}$.
\end{proposition}

\begin{proof}
Let $C_{\lambda,k}$ and $S_{\lambda,k}$ be the fundamental solutions with
data $(1,0)$ and $(0,1)$ at frequency $k\ge1$. Their Volterra equations and
Gronwall's inequality imply
$|C_{\lambda,k}|\le e^{|\lambda|Q/k}$ and
$|S_{\lambda,k}|\le e^{|\lambda|Q/k}/k$. Substitution back into those
equations bounds the differences from the free multipliers by
\[
 |C_{\lambda,k}-\cos k\eta|\le B/k,\qquad
 |S_{\lambda,k}-\sin(k\eta)/k|\le B/k^2,
 \quad B=|\lambda|Qe^{|\lambda|Q}.
\]
Acting on $\Lambda f$ rather than $f$, the first multiplier acquires another
$1/k$. The addition identity for an orthonormal harmonic basis is
$\sum_j|Y_{kj}(X)|^2=k^2/(2\pi^2)$, since the eigenspace has dimension
$k^2$ and its diagonal projector is rotation-invariant. Joint
Cauchy--Schwarz on the coefficients of $(\Lambda f,v)$ bounds the squared
evaluation norm by
\[
 \sum_{k=1}^\infty\frac{k^2}{2\pi^2}
       \left(\frac{B^2}{k^4}+\frac{B^2}{k^4}\right)
 =\frac{B^2}{6}.
\]
Removing the position-data term halves this bound. The argument concerns
smooth solutions and the bounded difference functional; it does not assign
point values to arbitrary individual finite-energy fields.
\end{proof}

Apply this to any rescaled difference $h_m=a(\phi_1-\phi_2)$ of two massive
solutions agreeing exactly on $\Is(b)$, with data $(f,v)$. Smooth extension to zero
implies that both data vanish on $B(Y,\tau)$. The massless solution $h_0$
with the same data vanishes on $\Is(b)$ by finite propagation. Thus
$h_0(x)=0$ for every $x\in F_0(b)$, and \eqref{eq:norm} gives
\begin{equation}\label{eq:hidden}
 |h_m(\eta,X)|\le
 \frac{m^2Q(\eta)e^{m^2Q(\eta)}}{\sqrt6}N(f,v),
 \qquad (\eta,X)\in F_0(b),
\end{equation}
where, for dust,
\[
 Q(\eta)=\frac{a_m^2}{4}
 \left(\frac{3\eta}{2}-2\sin\eta+\frac{\sin2\eta}{4}\right).
\]
The bound controls all such smooth alternatives with $N\le1$, not just the
witness in Proposition~\ref{prop:mass}. At fixed targets their rescaled
ambiguity is $O(m^2)$. For physical fields divide by $a(\eta)>0$; uniform
physical-field control requires staying away from the singular endpoints.
Without a norm bound any nonzero hidden solution can instead be scaled to
an arbitrarily large target difference. Neither convention is an observation
model supplied by the geometry.

\subsection{Errors in records are a different question}

Equation~\eqref{eq:hidden} retains exact agreement on the whole past while
changing the law. It is not a stability theorem for noisy records under a
fixed law. Even at a massless-fixed target, an aggregate field-error norm
need not control a point value. To see the regularity issue explicitly, fix
any interior target $(\eta_*,Y)$ and let $Z_k$ be the real unit-$L^2$ zonal
harmonic centered at $Y$, with $Z_k(Y)=k/\sqrt{2\pi^2}$. The smooth massless
solutions
\[
 \psi_K(\eta,X)=\frac1{\sqrt K}
  \sum_{k=K}^{2K-1}\frac{\cos(k(\eta-\eta_*))}{k}Z_k(X)
\]
have conserved conformal norm $N=1$ and target value
$\sqrt{K/(2\pi^2)}$, but
\[
 \|\psi_K\|_{L^2((0,2\pi)\times S^3)}^2
 =\frac\pi K\sum_{k=K}^{2K-1}\frac1{k^2}\le\frac\pi{K^2}.
\]
Their error on $\Is(b)$ is even smaller. They do not agree \emph{exactly}
there. This is unbounded point evaluation under a specified norm, already
present with records of the entire strip, not a causal-horizon or detector
failure. Every fixed spatial average with profile $\chi\in L^2(S^3)$ instead
satisfies $|\langle\chi,\psi_K(\eta)\rangle|\le\|\chi\|_2/K$.
No general past-record stability theorem for such averages is asserted.

A common finite spectral cutoff changes the exact question too: every
solution in that restricted massless class is a finite real-analytic sum,
so agreement on any nonempty open region entails agreement everywhere.
Its fixing set is all of $M$. Thus finite mode sums cannot establish exact
past agreement with a nonzero future difference in the unrestricted smooth
theory. These are classical regularity observations, not further priority
claims or results about the physical measurement process.
